% Figure: Euler hyperbola versus the Johnson parabola for an inelastic
% column. The Johnson (parabolic) formula caps the short-column end at
% the squash stress and meets the Euler hyperbola tangentially at the
% transition slenderness. A demand point marks the worked case study.
\begin{figure}[htbp]
\centering
\begin{tikzpicture}
\begin{axis}[
  width=12cm, height=7.5cm,
  xlabel={slenderness ratio $\lambda=KL/r$},
  ylabel={critical stress $\sigma_{\mathrm{cr}}$ (MPa)},
  xmin=0, xmax=180, ymin=0, ymax=360,
  axis lines=left,
  legend pos=north east, legend cell align=left,
  legend style={font=\scriptsize},
  tick label style={font=\scriptsize}, label style={font=\small},
  grid=major, grid style={enggray!18},
]
% Transition slenderness for sigma_y = 250, E = 200000 MPa:
% lambda_c = pi*sqrt(2E/sigma_y) = pi*sqrt(1600) = 125.7 (CRC convention)
% Johnson parabola: sigma_cr = sigma_y * (1 - lambda^2/(2 lambda_c^2)),
% with lambda_c = pi*sqrt(2E/sigma_y); plotted 0..125.7.
\addplot[engblue, very thick, domain=0:125.7, samples=80]
  {250*(1 - x^2/(2*125.7^2))};
\addlegendentry{Johnson parabola (inelastic)}
% Euler hyperbola sigma_cr = pi^2 E / lambda^2 for lambda > lambda_c.
\addplot[engred, very thick, domain=125.7:180, samples=80]
  {pi^2*200000/x^2};
\addlegendentry{Euler $\pi^2E/\lambda^2$}
% squash cap reference line
\addplot[enggray, dashed, thick] coordinates {(0,250) (180,250)};
\node[enggray, font=\scriptsize, anchor=south west] at (axis cs:2,252)
  {squash stress $\sigma_y=250$};
% transition slenderness marker
\addplot[enggray, dashed, thick] coordinates {(125.7,0) (125.7,125)};
\node[enggray, font=\scriptsize, anchor=west] at (axis cs:127,170)
  {$\lambda_c=\pi\sqrt{2E/\sigma_y}\approx 126$};
% demand point of the worked case study: lambda = 70, sigma about 211
\addplot[engamber, only marks, mark=*, mark size=2.6pt]
  coordinates {(70,211)};
\node[engamber, font=\scriptsize, anchor=west] at (axis cs:73,211)
  {case study $\lambda=70$};
\end{axis}
\end{tikzpicture}
\caption[Euler hyperbola and Johnson parabola for an inelastic
column]{The two branches of the column strength curve in the CRC
(Column Research Council) convention. Below the transition slenderness
$\lambda_c=\pi\sqrt{2E/\sigma_y}$ the inelastic Johnson parabola
governs, falling from the squash stress $\sigma_y$ at zero slenderness
and meeting the Euler hyperbola tangentially at $\lambda_c$, where the
Euler stress has dropped to exactly half the yield stress. Above
$\lambda_c$ the elastic Euler hyperbola governs. The amber point is the
strut of the case study at $\lambda=70$, comfortably inside the
inelastic range, where the parabola returns a critical stress well
above the elastic Euler value the unwary designer would read off the
hyperbola.}
\label{fig:vol03ch08:strut-euler-johnson}
\end{figure}
